% ==============================================================================
% chapters/ch03_architecture.tex - Chapter 3: System Architecture & Spec (~9 pp)
% ==============================================================================
\chapter{ארכיטקטורת המערכת ומפרט הנדסי}
\label{ch:architecture}

פרק זה מציג את הארכיטקטורה ההנדסית המלאה של מערכת \en{CivicShield AI}. התכנון מבוסס על עקרונות של הפרדת רשויות (\en{Separation of Concerns}), עמידות בכשלים (\en{Fault Tolerance}), אבטחה מובנית (\en{Security by Design}) ושקיפות מלאה הנדרשת במערכות בינה מלאכותית במגזר הציבורי. הפרק כולל את פירוט ארכיטקטורת חמש השכבות, מפרט רכזת הסוכנים (\en{Multi-Agent Hub}), טבלאות דרישות פונקציונליות (\en{FR}) ולא-פונקציונליות (\en{NFR}), וניתוח סיכונים ארכיטקטוניים.

% ------------------------------------------------------------------------------
\section{ארכיטקטורת חמש השכבות (\en{5-Layer Architecture})}
\label{sec:five-layer}

המערכת בנויה מחמש שכבות אופקיות מופרדות לוגית ורשתית, המלוות בשני עמודי תווך אנכיים חוצי-שכבות: \textbf{ממשל, אבטחה ופרטיות} (\en{Zero-Trust Governance}) ו-\textbf{ניטור, ביקורת ותצפיתיות} (\en{Observability \& Audit}). כל תקשורת בין שכבות מתבצעת באמצעות ממשקים מוגדרים היטב (\en{gRPC / mTLS}) בלבד, ללא גישה ישירה למסדי נתונים בין שכבות שאינן סמוכות.

\subsection{תיאור השכבות}

\begin{enumerate}[leftmargin=1.8em]
  \item \textbf{שכבה 1 -- שכבת איסוף ואינטגרציה (\en{Ingestion \& Sovereign Connectors}):}
  אחראית על קליטת מידע רציפה (\en{CDC / Streaming}) ובאצוות (\en{Batch}) ממאגרים ממשלתיים, מערכות \en{CRM}, תקנות חוק, פרוטוקולים ומסמכי \en{PDF/DOCX}. השכבה מבצעת ניקוי, פירוק ליחידות סמנטיות (\en{Chunking}), תיוג סיווג ביטחוני ומיסוך ראשוני של \en{PII} בטרם כניסת המידע למערכת.

  \item \textbf{שכבה 2 -- שכבת נתונים, ידע וזיכרון ארגוני (\en{Knowledge, Vector \& Audit Ledger}):}
  משלבת שלושה מנועי אחסון משלימים: (א) מאגר וקטורי היברידי (\en{Dense + BM25}) לשליפה מדויקת בעברית ובאנגלית; (ב) גרף ידע ארגוני (\en{Knowledge Graph}) למיפוי קשרי חוקים, סמכויות ונהלים; (ג) יומן ביקורת בלתי ניתן לשינוי (\en{Append-Only Cryptographic Audit Ledger}) הרושם כל שאילתה, מקור והחלטת סוכן.

  \item \textbf{שכבה 3 -- שכבת תזמור רב-סוכנים (\en{Multi-Agent Orchestration Hub}):}
  ליבת קבלת ההחלטות של \en{CivicShield AI}. מנהלת את מחזור החיים של הסוכנים המתמחים, מבצעת ניתוב דטרמיניסטי או סמנטי של משימות, אוכפת תקציבי השהיה וטוקנים, ומפעילה לולאות אימות הדדי בין סוכן יצרן לסוכן מבקר (\en{Critic Agent}).

  \item \textbf{שכבה 4 -- שכבת אבטחה, ציות ובקרת איכות (\en{Guardrails \& Policy Engine}):}
  פועלת כחומת אש קוגניטיבית דו-כיוונית: בכניסה היא חוסמת התקפות \en{Prompt Injection} ודליפת מידע רגיש; ביציאה היא מוודאת עיגון מלא במקורות (\en{Citation Grounding}), מניעת הזיות, עמידה ברגולציה והפניה לאישור אנושי (\en{HITL}) כאשר רמת הביטחון נמוכה מהסף המוגדר.

  \item \textbf{שכבה 5 -- שכבת ממשק, \en{API} והגשה (\en{Gov-Portal \& Integration Gateway}):}
  מספקת פורטל עבודה לעובדי מדינה, ממשקי \en{REST/GraphQL} לשילוב בפורטלים קיימים, מערכת התרעות למרכזי שליטה (\en{SIEM/SOC}), ולוח בקרה ניהולי להצגת מדדי \en{SLA}, דיוק וציות.
\end{enumerate}

\begin{figure}[htbp]
\centering
\begin{ltrblock}
\resizebox{\textwidth}{!}{%
\begin{tikzpicture}[
  font=\sffamily\footnotesize,
  layer/.style={
    rectangle, draw=CSNavy, thick, rounded corners=2pt,
    minimum width=11.4cm, text width=11.0cm, minimum height=1.22cm,
    inner sep=5pt, align=center, drop shadow,
    execute at begin node={\hyphenpenalty=10000\exhyphenpenalty=10000}
  },
  pillar/.style={
    rectangle, draw=CSGold, thick, rounded corners=2pt,
    fill=CSContractBg, minimum width=2.55cm, minimum height=8.15cm,
    align=center, text width=2.3cm,
    execute at begin node={\hyphenpenalty=10000\exhyphenpenalty=10000}
  },
  arr/.style={-{Stealth[length=2.5mm]}, thick, CSNavy}
]
  % 5 Horizontal Layers
  \node[layer, fill=CSNavy!12] (L5) at (0, 6.88)
    {\textbf{Layer 5: Gov-Portal UI, REST/gRPC Gateway \& SIEM}\\[2pt]
     {\scriptsize Citizen/Officer Portal $\cdot$ OAuth2/OIDC $\cdot$ Rate Limiting $\cdot$ Executive Dashboards}};

  \node[layer, fill=CSRiskBg, draw=CSRiskRed] (L4) at (0, 5.16)
    {\textbf{Layer 4: Zero-Trust Guardrails \& Policy Engine}\\[2pt]
     {\scriptsize Prompt-Injection Firewall $\cdot$ PII Redaction $\cdot$ Grounding Verifier $\cdot$ HITL Gate}};

  \node[layer, fill=CSInsightBg, draw=CSInsightTeal] (L3) at (0, 3.44)
    {\textbf{Layer 3: Multi-Agent Orchestration Hub}\\[2pt]
     {\scriptsize Router $\cdot$ Policy Agent $\cdot$ RAG Agent $\cdot$ Risk Scorer $\cdot$ Critic Agent $\cdot$ State Machine}};

  \node[layer, fill=CSIce, draw=CSSlate] (L2) at (0, 1.72)
    {\textbf{Layer 2: Knowledge Graph, Vector Store \& Audit Ledger}\\[2pt]
     {\scriptsize Hybrid Vector DB (HE/EN) $\cdot$ Legal Knowledge Graph $\cdot$ Immutable WORM Audit Log}};

  \node[layer, fill=CSNavy!8, draw=CSSteel] (L1) at (0, 0.0)
    {\textbf{Layer 1: Ingestion \& Sovereign Data Connectors}\\[2pt]
     {\scriptsize Gov ERP/CRM $\cdot$ Legal Gazettes $\cdot$ OCR/Chunking Pipeline $\cdot$ Classification Tagger}};

  % Vertical Cross-Cutting Pillars
  \node[pillar] (PL) at (-7.35, 3.44)
    {\textbf{\color{CSNavy}Zero-Trust\\Security}\\[6pt]
     mTLS 1.3\\[3pt]
     ABAC/RBAC\\[3pt]
     HSM/KMS\\[3pt]
     INCD Policy\\[3pt]
     Air-Gap Ready};

  \node[pillar] (PR) at (7.35, 3.44)
    {\textbf{\color{CSNavy}Observability\\\& Ops}\\[6pt]
     OpenTelemetry\\[3pt]
     Audit Trail\\[3pt]
     Drift Alerts\\[3pt]
     SLA Monitor\\[3pt]
     FinOps Quota};

  % Inter-layer bidirectional arrows
  \foreach \a/\b in {L1/L2, L2/L3, L3/L4, L4/L5} {
    \draw[arr] ([xshift=-1.6cm]\a.north) -- ([xshift=-1.6cm]\b.south);
    \draw[arr] ([xshift=1.6cm]\b.south) -- ([xshift=1.6cm]\a.north);
  }
\end{tikzpicture}%
}
\end{ltrblock}
\caption{ארכיטקטורת חמש השכבות של מערכת \en{CivicShield AI} ועמודי התווך חוצי-השכבות}
\label{fig:five-layer}
\end{figure}

\begin{insightbox}[עיקרון הניידות הריבונית (\en{Sovereign Portability})]
הפרדת שכבה 3 (תזמור) משכבה 2 (נתונים) באמצעות מתאמים סטנדרטיים מאפשרת החלפת מודל שפה (\en{LLM Swap}) או מעבר בין ענן ממשלתי (\en{Nimbus}) למרכז נתונים מאובטח (\en{On-Premise}) ללא שינוי בלוגיקה העסקית או במדיניות הציות.
\end{insightbox}

% ------------------------------------------------------------------------------
\section{מפרט רכזת הסוכנים הרב-מערכתית (\en{Multi-Agent Hub Spec})}
\label{sec:agent-hub}

רכזת הסוכנים (\en{Multi-Agent Hub}) מממשת ארכיטקטורת \en{Hub-and-Spoke} מונחית-אירועים (\en{Event-Driven}) המבוססת על תור הודעות אמין (\en{NATS JetStream / Kafka}) ומכונת מצבים דטרמיניסטית (\en{Deterministic Finite State Machine}). אף סוכן אינו רשאי לפנות ישירות לסוכן אחר ללא מעבר דרך הרכזת, וזאת כדי להבטיח אכיפת הרשאות, תיעוד מלא ומניעת לולאות אינסופיות.

\subsection{טקסונומיית הסוכנים ותחומי אחריות}

\begin{tabularx}{\textwidth}{@{} l l X @{}}
\toprule
\rowcolor{CSNavy}
\thead{מזהה סוכן} & \thead{תפקיד מערכתי} & \thead{אחריות ומגבלות גישה} \\
\midrule
\code{hub.router} & מנתב ומנהל זרימה & מפרק בקשות לתתי-משימות, מקצה תקציב זמן/טוקנים ומנהל \en{State}. \\
\code{agent.rag} & סוכן שליפה וידע & גישת קריאה בלבד לשכבה 2; מחזיר קטעי מקור עם \en{DocID} וציון רלוונטיות. \\
\code{agent.policy} & סוכן מדיניות וחוק & בודק תאימות מול גרף החוקים והנהלים; מזהה סתירות רגולטוריות. \\
\code{agent.risk} & סוכן הערכת סיכונים & מחשב ציון סיכון משוקלל ($R \in [0,1]$) ומפעיל הסלמה ל-\en{HITL} מעל סף. \\
\code{agent.critic} & סוכן ביקורת ואימות & מאמת כל טענה מול מקורות \code{agent.rag}; פוסל טענות לא מעוגנות. \\
\code{agent.synth} & סוכן ניסוח והגשה & מפיק תשובה מובנית בעברית תקנית עם הפניות מלאות ומטא-דאטה. \\
\bottomrule
\end{tabularx}

\subsection{טופולוגיית הרכזת וזרימת הבקרה}

\begin{figure}[htbp]
\centering
\begin{ltrblock}
\begin{tikzpicture}[
  font=\sffamily\footnotesize,
  hub/.style={
    rectangle, draw=CSNavy, very thick, rounded corners=4pt,
    fill=CSNavy, text=white, minimum width=4.2cm, minimum height=1.5cm,
    align=center, drop shadow
  },
  agent/.style={
    rectangle, draw=CSSlate, thick, rounded corners=2pt,
    fill=CSIce, minimum width=3.1cm, minimum height=1.0cm,
    align=center
  },
  critic/.style={
    rectangle, draw=CSRiskRed, thick, rounded corners=2pt,
    fill=CSRiskBg, minimum width=3.3cm, minimum height=1.0cm,
    align=center
  },
  bus/.style={-{Stealth[length=2.2mm]}, thick, CSSlate}
]
  % Central Hub
  \node[hub] (HUB) at (0, 0) {\textbf{Multi-Agent Hub}\\State Machine $\cdot$ Policy Gate\\DLQ \& Circuit Breaker};

  % Spokes
  \node[agent] (RAG) at (-4.8, 2.1) {\textbf{RAG Agent}\\Hybrid Search (L2)};
  \node[agent] (POL) at (0, 2.5) {\textbf{Policy Agent}\\Legal Graph Check};
  \node[agent] (RSK) at (4.8, 2.1) {\textbf{Risk Agent}\\Anomaly \& Score};

  \node[agent] (SYN) at (-4.5, -2.2) {\textbf{Synthesizer Agent}\\Hebrew Output + Citations};
  \node[critic] (CRT) at (4.5, -2.2) {\textbf{Critic Agent (Red-Check)}\\Grounding \& Hallucination Gate};

  % Connections
  \draw[bus, <->] (HUB.north west) -- (RAG.south east);
  \draw[bus, <->] (HUB.north) -- (POL.south);
  \draw[bus, <->] (HUB.north east) -- (RSK.south west);
  \draw[bus, <->] (HUB.south west) -- (SYN.north east);
  \draw[bus, <->] (HUB.south east) -- (CRT.north west);

  % Verification loop note
  \draw[bus, dashed, CSRiskRed] (SYN.east) -- node[below, text=CSRiskRed]{\scriptsize Mandatory Critic Pass} (CRT.west);
\end{tikzpicture}
\end{ltrblock}
\caption{טופולוגיית \en{Hub-and-Spoke} של רכזת הסוכנים ולולאת האימות ההכרחית}
\label{fig:agent-hub}
\end{figure}

\begin{contractbox}[חוזה מעטפת הודעה ו-\en{SLA} בין סוכנים (\en{Agent-to-Hub Contract})]
כל הודעה העוברת ברכזת חייבת להיות חתומה קריפטוגרפית (\en{Ed25519}) ולהכיל את השדות המחייבים הבאים:
\begin{ltrblock}
\ttfamily\footnotesize\color{CSNavy}
\{\\
\hspace*{1em}"msg\_id": "uuid-v7",\\
\hspace*{1em}"trace\_id": "otel-trace-128bit",\\
\hspace*{1em}"idempotency\_key": "sha256(session\_id + step\_seq)",\\
\hspace*{1em}"sender": "agent.rag",\\
\hspace*{1em}"recipient": "hub.router",\\
\hspace*{1em}"classification": "RESTRICTED\_GOV",\\
\hspace*{1em}"ttl\_ms": 600,\\
\hspace*{1em}"citations": ["DOC-IL-2026-881\#p4"],\\
\hspace*{1em}"confidence": 0.94,\\
\hspace*{1em}"signature": "base64(ed25519\_sig)"\\
\}
\end{ltrblock}
\textbf{חריגת \en{SLA}:} אם סוכן אינו מגיב בתוך \code{ttl\_ms}, הרכזת מפעילה \en{Circuit Breaker}, מעבירה את המשימה ל-\en{Fallback} דטרמיניסטי או ל-\en{HITL}, ורושמת אירוע ב-\en{DLQ}.
\end{contractbox}

% ------------------------------------------------------------------------------
\section{דרישות פונקציונליות (\en{Functional Requirements -- FR})}
\label{sec:fr}

טבלה \ref{tab:fr} מפרטת את הדרישות הפונקציונליות המחייבות למכרז. כל דרישה משויכת לרמת עדיפות (\en{Must / Should}) ולשיטת אימות רשמית בשלב מבחני הקבלה (\en{ATP}).

\begin{longtable}{@{} l p{3.1cm} p{6.8cm} l l @{}}
\caption{מטריצת דרישות פונקציונליות מחייבות (\en{FR-01} עד \en{FR-10})}
\label{tab:fr} \\
\toprule
\rowcolor{CSNavy}
\thead{מזהה} & \thead{תחום} & \thead{תיאור הדרישה הפונקציונלית} & \thead{עדיפות} & \thead{אימות} \\
\midrule
\endfirsthead
\toprule
\rowcolor{CSNavy}
\thead{מזהה} & \thead{תחום} & \thead{תיאור הדרישה הפונקציונלית} & \thead{עדיפות} & \thead{אימות} \\
\midrule
\endhead
\code{FR-01} & קליטת מסמכים & קליטה אוטומטית של \en{PDF}, \en{DOCX}, \en{HTML} ו-\en{JSON} עם \en{OCR} בעברית, ערבית ואנגלית. & \en{Must} & בדיקה \\
\code{FR-02} & חיפוש היברידי & שליפה סמנטית ומילונית משולבת (\en{Vector + BM25}) עם סינון לפי הרשאות \en{ABAC}. & \en{Must} & בדיקה \\
\code{FR-03} & תזמור סוכנים & ניתוב דינמי בין לפחות 6 סוכנים מתמחים באמצעות מכונת מצבים דטרמיניסטית. & \en{Must} & הדגמה \\
\code{FR-04} & עיגון ומראי מקום & כל טענה עובדתית בפלט תלווה בקישור ישיר למסמך המקור ולמספר עמוד/סעיף. & \en{Must} & בדיקה \\
\code{FR-05} & בקרת \en{Critic} & אימות עצמאי של הפלט מול מקורות ומדיניות בטרם הצגה למשתמש; חסימת הזיות. & \en{Must} & בדיקה \\
\code{FR-06} & הסלמה אנושית & העברה אוטומטית לתור \en{HITL} כאשר ציון הביטחון $< 0.85$ או זוהה סיכון משפטי. & \en{Must} & הדגמה \\
\code{FR-07} & מיסוך \en{PII} & זיהוי והסוואה אוטומטית של מספרי ת"ז, טלפונים וכתובות בלוגים ובפרומפטים. & \en{Must} & אבטחה \\
\code{FR-08} & יומן ביקורת & תיעוד קריפטוגרפי (\en{WORM}) של כל שאילתה, מקורות, תוצאות סוכנים ואישור אנושי. & \en{Must} & ביקורת \\
\code{FR-09} & עדכון מדיניות & טעינה חמה (\en{Hot-Reload}) של כללי ציות וחוקים חדשים ללא השבתת המערכת. & \en{Should} & הדגמה \\
\code{FR-10} & ממשק \en{SIEM} & הזרמת אירועי אבטחה וחריגות בזמן אמת בפורמט \en{CEF/Syslog} למרכז ה-\en{SOC}. & \en{Must} & אינטגרציה \\
\bottomrule
\end{longtable}

% ------------------------------------------------------------------------------
\section{דרישות לא-פונקציונליות (\en{Non-Functional Requirements -- NFR})}
\label{sec:nfr}

טבלה \ref{tab:nfr} מגדירה את מדדי האיכות, הביצועים, הזמינות והאבטחה המחייבים. איור \ref{fig:latency-budget} ממחיש את חלוקת תקציב ההשהיה (\en{Latency Budget}) מקצה לקצה.

\begin{longtable}{@{} l >{\raggedleft\arraybackslash\hyphenpenalty=10000\exhyphenpenalty=10000}p{3.4cm} p{6.6cm} l @{}}
\caption{מטריצת דרישות לא-פונקציונליות (\en{NFR-01} עד \en{NFR-10})}
\label{tab:nfr} \\
\toprule
\rowcolor{CSNavy}
\thead{מזהה} & \thead{קטגוריה} & \thead{מדד מחייב (\en{Target Metric})} & \thead{תקן / בסיס} \\
\midrule
\endfirsthead
\toprule
\rowcolor{CSNavy}
\thead{מזהה} & \thead{קטגוריה} & \thead{מדד מחייב (\en{Target Metric})} & \thead{תקן / בסיס} \\
\midrule
\endhead
\code{NFR-01} & השהיה (\en{Latency}) & זמן תגובה \en{P50} $\le 850\text{ms}$, \en{P95} $\le 1800\text{ms}$, \en{P99} $\le 2800\text{ms}$. & \en{Gov-SLA} \\
\code{NFR-02} & תפוקה (\en{Throughput}) & תמיכה ב-250 שאילתות בשנייה (\en{QPS}) בשגרה ועד 1,000 \en{QPS} בחירום. & מבחן עומס \\
\code{NFR-03} & זמינות (\en{Availability}) & זמינות של $99.95\%$ בתצורת \en{Active-Active} על פני שני אזורי זמינות. & \en{Tier-3+} \\
\code{NFR-04} & התאוששות מאסון & \en{RPO} $\le 60$ שניות, \en{RTO} $\le 15$ דקות. & \en{DRP} \\
\code{NFR-05} & דיוק והזיות & שיעור טענות לא מעוגנות $< 0.5\%$ במדגם ביקורת חודשי. & \en{AI-QA} \\
\code{NFR-06} & ריבונות מידע & $100\%$ מהנתונים והאינדקסים מאוחסנים פיזית בגבולות מדינת ישראל. & חוק הפרטיות \\
\code{NFR-07} & הצפנה & \en{TLS 1.3} בתעבורה, \en{AES-256-GCM} במנוחה עם ניהול מפתחות \en{HSM}. & \en{INCD} \\
\code{NFR-08} & בידוד דיירים & הפרדה לוגית והצפנית מלאה בין משרדי ממשלה (\en{Multi-Tenant Isolation}). & \en{ISO 27001} \\
\code{NFR-09} & תצפיתיות & כיסוי $100\%$ של \en{Distributed Tracing} בתקן \en{OpenTelemetry}. & \en{CNCF} \\
\code{NFR-10} & נגישות & עמידה מלאה בתקן ישראלי ת"י 5568 (\en{WCAG 2.1 AA}) לממשקי המשתמש. & חוק שוויון \\
\bottomrule
\end{longtable}

\begin{figure}[htbp]
\centering
\begin{ltrblock}
\begin{tikzpicture}[
  font=\sffamily\scriptsize,
  stage/.style={
    rectangle, draw=CSNavy, thick, minimum height=1.05cm,
    align=center, rounded corners=1.5pt
  },
  arr/.style={-{Stealth[length=2mm]}, thick, CSSlate}
]
  \node[stage, fill=CSIce, minimum width=2.0cm] (S1) at (0,0)
    {\textbf{Gateway \& WAF}\\80 ms};
  \node[stage, fill=CSRiskBg, minimum width=2.1cm] (S2) at (2.45,0)
    {\textbf{L4 Input Guard}\\120 ms};
  \node[stage, fill=CSInsightBg, minimum width=2.3cm] (S3) at (5.05,0)
    {\textbf{L2 Hybrid RAG}\\250 ms};
  \node[stage, fill=CSNavy!12, minimum width=2.5cm] (S4) at (7.85,0)
    {\textbf{L3 Agents + Synth}\\950 ms};
  \node[stage, fill=CSContractBg, minimum width=2.3cm] (S5) at (10.65,0)
    {\textbf{L4 Critic \& Sign}\\400 ms};

  \draw[arr] (S1) -- (S2);
  \draw[arr] (S2) -- (S3);
  \draw[arr] (S3) -- (S4);
  \draw[arr] (S4) -- (S5);

  % Total Budget Bracket
  \draw[decorate, decoration={brace, amplitude=6pt, mirror}, thick, CSGold]
    (S1.south west) -- (S5.south east)
    node[midway, below=8pt, text=CSNavy, font=\sffamily\footnotesize\bfseries]
    {Total End-to-End P95 Latency Budget $\le$ 1,800 ms};
\end{tikzpicture}
\end{ltrblock}
\caption{תקציב השהיה מקצה לקצה (\en{P95 Latency Budget}) לאורך צינור העיבוד}
\label{fig:latency-budget}
\end{figure}

% ------------------------------------------------------------------------------
\section{ניהול סיכונים ארכיטקטוניים ומוכנות לסקירת צוות אדום}
\label{sec:arch-risks}

במסגרת תכנון המערכת זוהו שלושה סיכונים ארכיטקטוניים מרכזיים, אשר עבורם הוטמעו בקרות מבניות כבר בשלב התכנון:

\begin{riskbox}[סיכון 1: הזרקת הנחיות (\en{Indirect Prompt Injection}) וחטיפת סוכן]
מסמך חיצוני שנקלט בשכבה 1 עלול להכיל הנחיות זדוניות המנסות להטות את \code{agent.synth} לחשוף מידע או לעקוף מדיניות.
\textbf{בקרת נגד:} הפרדה מבנית בין ערוץ ההנחיות (\en{System Policy}) לערוץ הנתונים (\en{Untrusted Context}), סריקת קלט בשכבה 4, ואימות בלתי תלוי של \code{agent.critic} שאינו חשוף לטקסט הגולמי של המשתמש.
\end{riskbox}

\begin{riskbox}[סיכון 2: מפל הזיות (\en{Hallucination Cascade}) בין סוכנים]
שגיאה מוקדמת של \code{agent.rag} או \code{agent.policy} עלולה להתעצם לאורך שרשרת הסוכנים.
\textbf{בקרת נגד:} חובת צירוף מזהי מקור (\code{citations}) בכל הודעת ביניים, חסימת טענות ללא עיגון, והסלמה אוטומטית ל-\en{HITL} בכל סתירה בין \code{agent.policy} ל-\code{agent.critic}.
\end{riskbox}

\begin{riskbox}[סיכון 3: נעילת ספק (\en{Vendor Lock-In}) וסיכון שרשרת אספקה]
תלות במודל שפה קנייני יחיד מסכנת את המשכיות השירות הממשלתי.
\textbf{בקרת נגד:} שכבת הפשטה אחידה (\en{Model Gateway}) התומכת הן במודלים פתוחים מאוחסנים מקומית (\en{Self-Hosted Weights}) והן בממשקים מאושרים בענן הריבוני.
\end{riskbox}

\begin{insightbox}[סטטוס מוכנות לסקירת צוות אדום (\en{Red-Team Readiness})]
ארכיטקטורת חמש השכבות, מפרט רכזת הסוכנים, חוזה ה-\en{SLA}, וטבלאות ה-\en{FR/NFR} הושלמו ומוכנים לבחינת חדירות, תרחישי קיצון והתקפות יריבות על ידי צוות \en{redteam} (פרק 5).
\end{insightbox}
